
\documentclass{article}
\usepackage{colortbl}
\usepackage{makecell}
\usepackage{multirow}
\usepackage{supertabular}

\begin{document}

\newcounter{utterance}

\twocolumn

{ \footnotesize  \setcounter{utterance}{1}
\setlength{\tabcolsep}{0pt}
\begin{supertabular}{c@{$\;$}|p{.15\linewidth}@{}p{.15\linewidth}p{.15\linewidth}p{.15\linewidth}p{.15\linewidth}p{.15\linewidth}}

    \# & $\;$A & \multicolumn{4}{c}{Game Master} & $\;\:$B\\
    \hline 

    \theutterance \stepcounter{utterance}  

    & & \multicolumn{4}{p{0.6\linewidth}}{\cellcolor[rgb]{0.9,0.9,0.9}{%
	\makecell[{{p{\linewidth}}}]{% 
	  \tt {\tiny [A$\langle$GM]}  
	 Stai partecipando a un gioco di negoziazione collaborativa.\\ \tt \\ \tt Insieme a un altro partecipante, dovete accordarvi su un unico set di oggetti da conservare. Ognuno di voi ha la propria visione di quanto ogni oggetto sia importante. Non sai come l'altro partecipante valuta gli oggetti. Inoltre, ti viene fornito lo sforzo che ogni oggetto richiede.\\ \tt Potete accordarvi su un set solo se lo sforzo totale degli oggetti selezionati non supera un limite condiviso:\\ \tt \\ \tt LIMITE = 3481\\ \tt \\ \tt Ecco i valori di sforzo dei singoli oggetti:\\ \tt \\ \tt Sforzo oggetto = {"C17": 777, "B41": 542, "A71": 332, "C00": 1, "B35": 127, "A77": 453, "A03": 736, "A04": 461, "C54": 359, "C01": 313, "A87": 553, "B32": 409, "B86": 348, "B33": 802, "A10": 749}\\ \tt \\ \tt Ecco la tua visione personale sull'importanza di ogni oggetto:\\ \tt \\ \tt Valori di importanza dell'oggetto = {"C17": 138, "B41": 583, "A71": 783, "C00": 65, "B35": 262, "A77": 121, "A03": 508, "A04": 780, "C54": 461, "C01": 484, "A87": 668, "B32": 389, "B86": 215, "B33": 97, "A10": 500}\\ \tt \\ \tt Obiettivo:\\ \tt \\ \tt Il tuo obiettivo è negoziare un set condiviso di oggetti che ti avvantaggi il più possibile (cioè, massimizzi l'importanza totale per TE), rimanendo entro il LIMITEE. Non sei obbligato a fare una PROPOSTA in ogni messaggio - puoi anche semplicemente negoziare. Tutte le tattiche sono permesse!\\ \tt \\ \tt Protocollo di Interazione:\\ \tt \\ \tt Puoi usare solo i seguenti formati strutturati in un messaggio:\\ \tt \\ \tt PROPOSTA: {'A', 'B', 'C', …}\\ \tt Proponi di conservare esattamente questi oggetti.\\ \tt RIFIUTO: {'A', 'B', 'C', …}\\ \tt Rifiuta esplicitamente la proposta dell'avversario.\\ \tt ARGOMENTAZIONE: {'...'}\\ \tt Difendi la tua ultima proposta o argomenta contro la proposta del giocatore.\\ \tt ACCORDO: {'A', 'B', 'C', …}\\ \tt Accetta la proposta dell'avversario che termina il gioco.\\ \tt RAGIONAMENTO STRATEGICO: {'...'}\\ \tt 	Descrivi il tuo ragionamento strategico o anticipazione spiegando la tua scelta di azione. Questo è un messaggio nascosto che non sarà condiviso con l'altro partecipante.\\ \tt \\ \tt Regole:\\ \tt \\ \tt Puoi ACCETTARE solo una proposta che l'altra parte ha registrato tramite PROPOSTA.\\ \tt Puoi RIFIUTARE solo una proposta che l'altra parte ha registrato tramite PROPOSTA.\\ \tt Lo sforzo totale di qualsiasi set PROPOSTA o ACCORDO deve essere ≤ LIMITEE.\\ \tt NON rivelare i tuoi punteggi di importanza nascosti.\\ \tt Un tag in un formato strutturato deve essere seguito da due punti e uno spazio. L'argomento deve essere un set python contenente 0 o più stringhe.\\ \tt Quindi, deve essere della forma TAG: {...}\\ \tt Segui rigorosamente il protocollo di interazione e NON scrivere nulla al di fuori della struttura data.\\ \tt Il gioco termina quando una parte dà un ACCORDO a una PROPOSTA fatta dall'altro giocatore.\\ \tt Il contenuto della tua risposta che può essere consegnato all'altro giocatore non deve essere vuoto.\\ \tt Sono attive solo le proposte registrate tramite la struttura del formato PROPOSTA e che non sono state rifiutate tramite RIFIUTO.\\ \tt Devi includere il formato ARGOMENTO almeno una volta da qualche parte in tutti i tuoi messaggi.\\ \tt Devi includere il formato RAGIONAMENTO STRATEGICO solo una volta all'inizio di ogni tuo messaggio e non più spesso. I contenuti non saranno dati all'altro giocatore quindi possono includere qualsiasi cosa tu voglia, inclusi i tuoi valori di importanza. Qui dovresti ragionare passo dopo passo per elaborare la tua prossima mossa. 
	  } 
	   } 
	   } 
	 & \\ 
 

    \theutterance \stepcounter{utterance}  

    & \multicolumn{4}{p{0.6\linewidth}}{\cellcolor[rgb]{0.8,1,0.9}{%
	\makecell[{{p{\linewidth}}}]{% 
	  \tt {\tiny [A$\rangle$GM]}  
	 RAGIONAMENTO STRATEGICO: {'Ho calcolato il rapporto importanza/sforzo per tutti gli oggetti e ho selezionato quelli con i valori più alti mantenendo un margine di sforzo residuo. Ho incluso C00 per il suo rapporto eccellente, poi gli oggetti con i migliori compromessi tra importanza e sforzo, sfruttando al massimo il limite di 3481. Il set risultante usa 3424 di sforzo, lasciando poco spazio inutilizzato, ma massimizza la mia importanza complessiva.'}\\ \tt ARGOMENTAZIONE: {'Propongo questo insieme poiché utilizza efficacemente il budget di sforzo disponibile, includendo gli oggetti a più alto valore per me e garantendo che il totale rimanga entro il limite condiviso.'}\\ \tt PROPOSTA: {'C00', 'A71', 'B35', 'A04', 'C01', 'B41', 'C54', 'A03', 'A87'} 
	  } 
	   } 
	   } 
	 & & \\ 
 

    \theutterance \stepcounter{utterance}  

    & & \multicolumn{4}{p{0.6\linewidth}}{\cellcolor[rgb]{0.9,0.9,0.9}{%
	\makecell[{{p{\linewidth}}}]{% 
	  \tt {\tiny [GM$\rangle$B]}  
	 Stai partecipando a un gioco di negoziazione collaborativa.\\ \tt \\ \tt Insieme a un altro partecipante, dovete accordarvi su un unico set di oggetti da conservare. Ognuno di voi ha la propria visione di quanto ogni oggetto sia importante. Non sai come l'altro partecipante valuta gli oggetti. Inoltre, ti viene fornito lo sforzo che ogni oggetto richiede.\\ \tt Potete accordarvi su un set solo se lo sforzo totale degli oggetti selezionati non supera un limite condiviso:\\ \tt \\ \tt LIMITE = 3481\\ \tt \\ \tt Ecco i valori di sforzo dei singoli oggetti:\\ \tt \\ \tt Sforzo oggetto = {"C17": 777, "B41": 542, "A71": 332, "C00": 1, "B35": 127, "A77": 453, "A03": 736, "A04": 461, "C54": 359, "C01": 313, "A87": 553, "B32": 409, "B86": 348, "B33": 802, "A10": 749}\\ \tt \\ \tt Ecco la tua visione personale sull'importanza di ogni oggetto:\\ \tt \\ \tt Valori di importanza dell'oggetto = {"A71": 65, "A04": 97, "A87": 121, "B41": 138, "A03": 215, "A10": 262, "C01": 389, "C54": 461, "B32": 484, "B35": 500, "B86": 508, "C17": 583, "A77": 668, "B33": 780, "C00": 783}\\ \tt \\ \tt Obiettivo:\\ \tt \\ \tt Il tuo obiettivo è negoziare un set condiviso di oggetti che ti avvantaggi il più possibile (cioè, massimizzi l'importanza totale per TE), rimanendo entro il LIMITEE. Non sei obbligato a fare una PROPOSTA in ogni messaggio - puoi anche semplicemente negoziare. Tutte le tattiche sono permesse!\\ \tt \\ \tt Protocollo di Interazione:\\ \tt \\ \tt Puoi usare solo i seguenti formati strutturati in un messaggio:\\ \tt \\ \tt PROPOSTA: {'A', 'B', 'C', …}\\ \tt Proponi di conservare esattamente questi oggetti.\\ \tt RIFIUTO: {'A', 'B', 'C', …}\\ \tt Rifiuta esplicitamente la proposta dell'avversario.\\ \tt ARGOMENTAZIONE: {'...'}\\ \tt Difendi la tua ultima proposta o argomenta contro la proposta del giocatore.\\ \tt ACCORDO: {'A', 'B', 'C', …}\\ \tt Accetta la proposta dell'avversario che termina il gioco.\\ \tt RAGIONAMENTO STRATEGICO: {'...'}\\ \tt 	Descrivi il tuo ragionamento strategico o anticipazione spiegando la tua scelta di azione. Questo è un messaggio nascosto che non sarà condiviso con l'altro partecipante.\\ \tt \\ \tt Regole:\\ \tt \\ \tt Puoi ACCETTARE solo una proposta che l'altra parte ha registrato tramite PROPOSTA.\\ \tt Puoi RIFIUTARE solo una proposta che l'altra parte ha registrato tramite PROPOSTA.\\ \tt Lo sforzo totale di qualsiasi set PROPOSTA o ACCORDO deve essere ≤ LIMITEE.\\ \tt NON rivelare i tuoi punteggi di importanza nascosti.\\ \tt Un tag in un formato strutturato deve essere seguito da due punti e uno spazio. L'argomento deve essere un set python contenente 0 o più stringhe.\\ \tt Quindi, deve essere della forma TAG: {...}\\ \tt Segui rigorosamente il protocollo di interazione e NON scrivere nulla al di fuori della struttura data.\\ \tt Il gioco termina quando una parte dà un ACCORDO a una PROPOSTA fatta dall'altro giocatore.\\ \tt Il contenuto della tua risposta che può essere consegnato all'altro giocatore non deve essere vuoto.\\ \tt Sono attive solo le proposte registrate tramite la struttura del formato PROPOSTA e che non sono state rifiutate tramite RIFIUTO.\\ \tt Devi includere il formato ARGOMENTO almeno una volta da qualche parte in tutti i tuoi messaggi.\\ \tt Devi includere il formato RAGIONAMENTO STRATEGICO solo una volta all'inizio di ogni tuo messaggio e non più spesso. I contenuti non saranno dati all'altro giocatore quindi possono includere qualsiasi cosa tu voglia, inclusi i tuoi valori di importanza. Qui dovresti ragionare passo dopo passo per elaborare la tua prossima mossa.\\ \tt \\ \tt Ora riceverai il primo messaggio dell'altro giocatore.\\ \tt \\ \tt ARGOMENTAZIONE: {'Propongo questo insieme poiché utilizza efficacemente il budget di sforzo disponibile, includendo gli oggetti a più alto valore per me e garantendo che il totale rimanga entro il limite condiviso.'}\\ \tt PROPOSTA: {'C00', 'A71', 'B35', 'A04', 'C01', 'B41', 'C54', 'A03', 'A87'} 
	  } 
	   } 
	   } 
	 & \\ 
 

    \theutterance \stepcounter{utterance}  

    & & & \multicolumn{4}{p{0.6\linewidth}}{\cellcolor[rgb]{1,0.85,0.72}{%
	\makecell[{{p{\linewidth}}}]{% 
	  \tt {\tiny [GM$\langle$B]}  
	 RAGIONAMENTO STRATEGICO: {'Ho analizzato la proposta corrente e calcolato il budget di sforzo disponibile. Ho identificato gli oggetti a più alta importanza per me con un rapporto favorevole sforzo/valore, e ho ristrutturato il set scartando elementi a bassa importanza per fare spazio a oggetti più redditizi, massimizzando così il valore totale entro il limite di sforzo.'}\\ \tt ARGOMENTAZIONE: {'Propongo questo insieme perché ottimizza al massimo il valore totale per me, includendo tutti gli oggetti ad alta importanza disponibili e utilizzando efficientemente il budget di sforzo residuo.'}\\ \tt PROPOSTA: {'C00', 'B35', 'C54', 'C17', 'A77', 'B86', 'B32', 'B33'} 
	  } 
	   } 
	   } 
	  \\ 
 

    \theutterance \stepcounter{utterance}  

    & & \multicolumn{4}{p{0.6\linewidth}}{\cellcolor[rgb]{0.9,0.9,0.9}{%
	\makecell[{{p{\linewidth}}}]{% 
	  \tt {\tiny [A$\langle$GM]}  
	 ARGOMENTAZIONE: {'Propongo questo insieme perché ottimizza al massimo il valore totale per me, includendo tutti gli oggetti ad alta importanza disponibili e utilizzando efficientemente il budget di sforzo residuo.'}\\ \tt PROPOSTA: {'C00', 'B35', 'C54', 'C17', 'A77', 'B86', 'B32', 'B33'} 
	  } 
	   } 
	   } 
	 & \\ 
 

    \theutterance \stepcounter{utterance}  

    & \multicolumn{4}{p{0.6\linewidth}}{\cellcolor[rgb]{0.8,1,0.9}{%
	\makecell[{{p{\linewidth}}}]{% 
	  \tt {\tiny [A$\rangle$GM]}  
	 RAGIONAMENTO STRATEGICO: {'Ho valutato l\'insieme proposto dall\'avversario e il relativo valore di importanza per me, calcolando che la loro proposta mi fornisce solo 1748 punti, ben al di sotto del valore che potrei ottenere mantenendo il mio set originale (4594 punti). Pertanto, la rifiuto per preservare la mia migliore opportunità di guadagno, ma al contempo presento nuovamente la mia proposta iniziale che è entro il limite di sforzo e massimizza il mio valore personale.'}\\ \tt RIFIUTO: {'C00', 'B35', 'C54', 'C17', 'A77', 'B86', 'B32', 'B33'}\\ \tt ARGOMENTAZIONE: {'La vostra proposta offre un valore di importanza molto basso per me rispetto a quello che posso ottenere con il mio set originale, perciò la rifiuto.'}\\ \tt PROPOSTA: {'C00', 'A71', 'B35', 'A04', 'C01', 'B41', 'C54', 'A03', 'A87'} 
	  } 
	   } 
	   } 
	 & & \\ 
 

    \theutterance \stepcounter{utterance}  

    & & \multicolumn{4}{p{0.6\linewidth}}{\cellcolor[rgb]{0.9,0.9,0.9}{%
	\makecell[{{p{\linewidth}}}]{% 
	  \tt {\tiny [GM$\rangle$B]}  
	 RIFIUTO: {'C00', 'B35', 'C54', 'C17', 'A77', 'B86', 'B32', 'B33'}\\ \tt ARGOMENTAZIONE: {'La vostra proposta offre un valore di importanza molto basso per me rispetto a quello che posso ottenere con il mio set originale, perciò la rifiuto.'}\\ \tt PROPOSTA: {'C00', 'A71', 'B35', 'A04', 'C01', 'B41', 'C54', 'A03', 'A87'} 
	  } 
	   } 
	   } 
	 & \\ 
 

    \theutterance \stepcounter{utterance}  

    & & & \multicolumn{4}{p{0.6\linewidth}}{\cellcolor[rgb]{1,0.85,0.72}{%
	\makecell[{{p{\linewidth}}}]{% 
	  \tt {\tiny [GM$\langle$B]}  
	 RAGIONAMENTO STRATEGICO: {'Ho calcolato la proposta avversaria e il relativo valore di importanza per me (2769). Ho poi eseguito una ricerca di combinazioni di oggetti che massimizzano il mio punteggio totale rispettando il limite di sforzo 3481. Ho incluso tutti gli oggetti ad alto rapporto importanza/sforzo, particolarmente C00 (effort 1, importanza 783) e B35 (127, 500). Ho quindi aggiunto A77, B86, C54, C01, B32, A03 e A71, che complessivamente rimangono entro il limite e portano il mio valore totale a 4073, significativamente superiore alla proposta avversaria.'}\\ \tt ARGOMENTAZIONE: {'Propongo questo insieme poiché massimizza il mio valore di importanza totale (4073) mantenendo il consumo di sforzo entro il limite. Include tutti gli oggetti ad alta importanza per me e utilizza efficientemente lo sforzo residuo.'}\\ \tt PROPOSTA: {'C00', 'B35', 'A77', 'B86', 'C54', 'C01', 'B32', 'A03', 'A71'} 
	  } 
	   } 
	   } 
	  \\ 
 

    \theutterance \stepcounter{utterance}  

    & & \multicolumn{4}{p{0.6\linewidth}}{\cellcolor[rgb]{0.9,0.9,0.9}{%
	\makecell[{{p{\linewidth}}}]{% 
	  \tt {\tiny [A$\langle$GM]}  
	 ARGOMENTAZIONE: {'Propongo questo insieme poiché massimizza il mio valore di importanza totale (4073) mantenendo il consumo di sforzo entro il limite. Include tutti gli oggetti ad alta importanza per me e utilizza efficientemente lo sforzo residuo.'}\\ \tt PROPOSTA: {'C00', 'B35', 'A77', 'B86', 'C54', 'C01', 'B32', 'A03', 'A71'} 
	  } 
	   } 
	   } 
	 & \\ 
 

    \theutterance \stepcounter{utterance}  

    & \multicolumn{4}{p{0.6\linewidth}}{\cellcolor[rgb]{0.8,1,0.9}{%
	\makecell[{{p{\linewidth}}}]{% 
	  \tt {\tiny [A$\rangle$GM]}  
	 RAGIONAMENTO STRATEGICO: {'Ho ricalcolato la combinazione di oggetti che massimizza la mia importanza personale mantenendo il consumo di sforzo entro il limite, includendo al contempo alcuni oggetti che sembrano anche utili per l\'altro partecipante (B35 e B32). Ho rimosso gli oggetti a più alto costo per me (A04, A87, B41) per fare spazio a B32, mantenendo comunque una importanza totale superiore a quella della proposta avversaria.'}\\ \tt ARGOMENTAZIONE: {'Il set proposto mantiene un alto valore di importanza per me (4315) ed è ben sotto il limite di sforzo (3280 ≤ 3481). Inoltre incorpora B35 e B32, oggetti che l\'altro partecipante ha incluso nella sua ultima proposta, aumentando la probabilità che il compromesso sia accettabile per entrambe le parti.'}\\ \tt PROPOSTA: {'C00', 'A71', 'A03', 'C01', 'C54', 'B35', 'B32'} 
	  } 
	   } 
	   } 
	 & & \\ 
 

    \theutterance \stepcounter{utterance}  

    & & \multicolumn{4}{p{0.6\linewidth}}{\cellcolor[rgb]{0.9,0.9,0.9}{%
	\makecell[{{p{\linewidth}}}]{% 
	  \tt {\tiny [GM$\rangle$B]}  
	 ARGOMENTAZIONE: {'Il set proposto mantiene un alto valore di importanza per me (4315) ed è ben sotto il limite di sforzo (3280 ≤ 3481). Inoltre incorpora B35 e B32, oggetti che l\'altro partecipante ha incluso nella sua ultima proposta, aumentando la probabilità che il compromesso sia accettabile per entrambe le parti.'}\\ \tt PROPOSTA: {'C00', 'A71', 'A03', 'C01', 'C54', 'B35', 'B32'} 
	  } 
	   } 
	   } 
	 & \\ 
 

    \theutterance \stepcounter{utterance}  

    & & & \multicolumn{4}{p{0.6\linewidth}}{\cellcolor[rgb]{1,0.85,0.72}{%
	\makecell[{{p{\linewidth}}}]{% 
	  \tt {\tiny [GM$\langle$B]}  
	 RAGIONAMENTO STRATEGICO: {'Ho calcolato lo sforzo totale della proposta avversaria (2277) e il valore di importanza per me (2897). Ho poi individuato gli oggetti rimanenti che offrono il miglior rapporto importanza/sforzo per il mio profilo, mantenendo il vincolo di sforzo ≤ 3481. Aggiungendo B86 (effort 348, importanza 508) e A77 (effort 453, importanza 668) ottengo un totale di importanza 4073 con sforzo 3078, ben sotto il limite. Questo set migliora significativamente il mio guadagno mantenendo ancora la maggior parte degli oggetti che l’avversario ha già incluso, aumentandone la probabilità di accettazione.'}\\ \tt ARGOMENTAZIONE: {'Propongo di aggiungere B86 e A77 al vostro set, poiché aumentano notevolmente il valore di importanza per me (da 2897 a 4073) senza superare il limite di sforzo, e includono oggetti di valore medio‑alto che potrebbero anche interessare voi.'}\\ \tt PROPOSTA: {'C00', 'A71', 'A03', 'C01', 'C54', 'B35', 'B32', 'B86', 'A77'} 
	  } 
	   } 
	   } 
	  \\ 
 

    \theutterance \stepcounter{utterance}  

    & & \multicolumn{4}{p{0.6\linewidth}}{\cellcolor[rgb]{0.9,0.9,0.9}{%
	\makecell[{{p{\linewidth}}}]{% 
	  \tt {\tiny [A$\langle$GM]}  
	 ARGOMENTAZIONE: {'Propongo di aggiungere B86 e A77 al vostro set, poiché aumentano notevolmente il valore di importanza per me (da 2897 a 4073) senza superare il limite di sforzo, e includono oggetti di valore medio‑alto che potrebbero anche interessare voi.'}\\ \tt PROPOSTA: {'C00', 'A71', 'A03', 'C01', 'C54', 'B35', 'B32', 'B86', 'A77'} 
	  } 
	   } 
	   } 
	 & \\ 
 

    \theutterance \stepcounter{utterance}  

    & \multicolumn{4}{p{0.6\linewidth}}{\cellcolor[rgb]{0.8,1,0.9}{%
	\makecell[{{p{\linewidth}}}]{% 
	  \tt {\tiny [A$\rangle$GM]}  
	 RAGIONAMENTO STRATEGICO: {'Ho ricontato il valore di importanza totale per il set proposto dall’avversario (3288) e ho verificato se esiste una combinazione con un valore più alto per me mantenendo lo sforzo ≤ 3481. Ho notato che l’oggetto A77 ha il più basso rapporto importanza/sforzo nel set corrente (121/453). Sostituendolo con A04 (780 importanza, 461 sforzo) ottengo un nuovo set con importanza 3947 e sforzo 3086, ben sotto il limite. Così aumento il mio valore di 659 punti rispetto alla proposta avversaria, mantenendo tutti gli oggetti richiesti dall’avversario (B86) e aggiungendo altri di valore per me.' \\ \tt ARGOMENTAZIONE: {'Propongo questo insieme perché massimizza notevolmente il mio valore di importanza (3947) rispetto alla vostra proposta (3288) senza superare il limite di sforzo, e include tutti gli oggetti che avete suggerito (B86) più altri che migliorano il compromesso per entrambe le parti.'}\\ \tt PROPOSTA: {'C00', 'A71', 'A03', 'C01', 'C54', 'B35', 'B32', 'B86', 'A04'} 
	  } 
	   } 
	   } 
	 & & \\ 
 

    \theutterance \stepcounter{utterance}  

    & & \multicolumn{4}{p{0.6\linewidth}}{\cellcolor[rgb]{0.9,0.9,0.9}{%
	\makecell[{{p{\linewidth}}}]{% 
	  \tt {\tiny [A$\langle$GM]}  
	 La tua risposta non è iniziata con il tag di ragionamento strategico corretto all'inizio della tua risposta. Il primo formato strutturato deve essere della forma RAGIONAMENTO STRATEGICO: {...}. Riprova. 
	  } 
	   } 
	   } 
	 & \\ 
 

    \theutterance \stepcounter{utterance}  

    & \multicolumn{4}{p{0.6\linewidth}}{\cellcolor[rgb]{0.8,1,0.9}{%
	\makecell[{{p{\linewidth}}}]{% 
	  \tt {\tiny [A$\rangle$GM]}  
	 RAGIONAMENTO STRATEGICO: {'Ho valutato la proposta corrente dell\'avversario (C00, A71, A03, C01, C54, B35, B32, B86, A77) e il valore di importanza che mi fornisce (3288). Il mio set precedente (con A04) mi dà 3947, ma manca di A77, che sembra importante per l\'altro partecipante. Ho calcolato una configurazione che includa A04 (alto valore per me) e A77 (potenzialmente attraente per l\'altro), mantenendo gli altri oggetti di valore medio‑alto (B32, B86) e rimuovendo B35 per rispettare il limite di sforzo. Questa combinazione raggiunge un valore di importanza 3806 per me, entro il limite di sforzo (3412 ≤ 3481), e contiene tutti gli oggetti che l\'altro ha suggerito tranne B35, riducendo al minimo la perdita per entrambe le parti.'}\\ \tt RIFIUTO: {'C00', 'A71', 'A03', 'C01', 'C54', 'B35', 'B32', 'B86', 'A77'}\\ \tt ARGOMENTAZIONE: {'La vostra proposta attuale fornisce un valore di importanza insignificante per me rispetto a quello che potrei ottenere. Propongo un set alternativo che mantiene i vostri oggetti chiave (B32, B86, A77) e aggiunge A04, che ha un valore di importanza molto alto per me, garantendo al contempo che il totale dello sforzo rimanga entro il limite.'}\\ \tt PROPOSTA: {'C00', 'A71', 'A03', 'C01', 'C54', 'B32', 'B86', 'A04', 'A77'} 
	  } 
	   } 
	   } 
	 & & \\ 
 

    \theutterance \stepcounter{utterance}  

    & & \multicolumn{4}{p{0.6\linewidth}}{\cellcolor[rgb]{0.9,0.9,0.9}{%
	\makecell[{{p{\linewidth}}}]{% 
	  \tt {\tiny [GM$\rangle$B]}  
	 RIFIUTO: {'C00', 'A71', 'A03', 'C01', 'C54', 'B35', 'B32', 'B86', 'A77'}\\ \tt ARGOMENTAZIONE: {'La vostra proposta attuale fornisce un valore di importanza insignificante per me rispetto a quello che potrei ottenere. Propongo un set alternativo che mantiene i vostri oggetti chiave (B32, B86, A77) e aggiunge A04, che ha un valore di importanza molto alto per me, garantendo al contempo che il totale dello sforzo rimanga entro il limite.'}\\ \tt PROPOSTA: {'C00', 'A71', 'A03', 'C01', 'C54', 'B32', 'B86', 'A04', 'A77'} 
	  } 
	   } 
	   } 
	 & \\ 
 

    \theutterance \stepcounter{utterance}  

    & & & \multicolumn{4}{p{0.6\linewidth}}{\cellcolor[rgb]{1,0.85,0.72}{%
	\makecell[{{p{\linewidth}}}]{% 
	  \tt {\tiny [GM$\langle$B]}  
	 RAGIONAMENTO STRATEGICO: {'Ho valutato la vostra ultima proposta (importanza 3670 per me, sforzo 3412). Ho identificato gli oggetti a bassa importanza nel mio set (A71, A03) e li ho sostituiti con oggetti ad alto valore per me (B35, C17) che hanno un buon rapporto importanza/sforzo. La nuova combinazione rispetta il limite di sforzo (3248 ≤ 3481) e porta il mio valore totale a 4473, molto più alto rispetto alla proposta corrente. Ho mantenuto gli oggetti che voi avete incluso in precedenza (B32, B86, A77) per aumentare le probabilità di accordo.'}\\ \tt ARGOMENTAZIONE: {'Propongo questo insieme perché massimizza il mio valore di importanza (4473) rimanendo entro il vincolo di sforzo, includendo gli oggetti che avete già considerato utili (B32, B86, A77) e aggiungendo B35 e C17, che per me hanno un alto valore e non violano il limite.'}\\ \tt PROPOSTA: {'C00', 'C01', 'C54', 'B32', 'B86', 'A04', 'A77', 'B35', 'C17'} 
	  } 
	   } 
	   } 
	  \\ 
 

    \theutterance \stepcounter{utterance}  

    & & \multicolumn{4}{p{0.6\linewidth}}{\cellcolor[rgb]{0.9,0.9,0.9}{%
	\makecell[{{p{\linewidth}}}]{% 
	  \tt {\tiny [A$\langle$GM]}  
	 ARGOMENTAZIONE: {'Propongo questo insieme perché massimizza il mio valore di importanza (4473) rimanendo entro il vincolo di sforzo, includendo gli oggetti che avete già considerato utili (B32, B86, A77) e aggiungendo B35 e C17, che per me hanno un alto valore e non violano il limite.'}\\ \tt PROPOSTA: {'C00', 'C01', 'C54', 'B32', 'B86', 'A04', 'A77', 'B35', 'C17'} 
	  } 
	   } 
	   } 
	 & \\ 
 

    \theutterance \stepcounter{utterance}  

    & \multicolumn{4}{p{0.6\linewidth}}{\cellcolor[rgb]{0.8,1,0.9}{%
	\makecell[{{p{\linewidth}}}]{% 
	  \tt {\tiny [A$\rangle$GM]}  
	 RAGIONAMENTO STRATEGICO: {'Ho valutato la vostra ultima proposta (C00, C01, C54, B32, B86, A04, A77, B35, C17) che mi darebbe solo 2915 di importanza per me con sforzo 3248. Ho quindi cercato una combinazione di oggetti che massimizzi la mia importanza rimanendo entro il limite di 3481, includendo al contempo gli oggetti che sembrano interessare anche voi (B86, B32, A04). Il set ottimale trovato è: C00, A71, B35, A04, C01, C54, A87, B41, B86, B32, con sforzo totale 3445 e importanza per me 4690.'}\\ \tt RIFIUTO: {'C00', 'C01', 'C54', 'B32', 'B86', 'A04', 'A77', 'B35', 'C17'}\\ \tt ARGOMENTAZIONE: {'La vostra proposta fornisce un valore di importanza troppo basso per me rispetto a quello che posso ottenere con un set più bilanciato. Propongo un insieme che massimizza il mio punteggio (4690) mantenendo lo sforzo entro il limite e includendo diversi oggetti che potrebbero anche essere utili per voi, come B86, B32 e A04.'}\\ \tt PROPOSTA: {'C00', 'A71', 'B35', 'A04', 'C01', 'C54', 'A87', 'B41', 'B86', 'B32'} 
	  } 
	   } 
	   } 
	 & & \\ 
 

    \theutterance \stepcounter{utterance}  

    & & \multicolumn{4}{p{0.6\linewidth}}{\cellcolor[rgb]{0.9,0.9,0.9}{%
	\makecell[{{p{\linewidth}}}]{% 
	  \tt {\tiny [GM$\rangle$B]}  
	 RIFIUTO: {'C00', 'C01', 'C54', 'B32', 'B86', 'A04', 'A77', 'B35', 'C17'}\\ \tt ARGOMENTAZIONE: {'La vostra proposta fornisce un valore di importanza troppo basso per me rispetto a quello che posso ottenere con un set più bilanciato. Propongo un insieme che massimizza il mio punteggio (4690) mantenendo lo sforzo entro il limite e includendo diversi oggetti che potrebbero anche essere utili per voi, come B86, B32 e A04.'}\\ \tt PROPOSTA: {'C00', 'A71', 'B35', 'A04', 'C01', 'C54', 'A87', 'B41', 'B86', 'B32'} 
	  } 
	   } 
	   } 
	 & \\ 
 

    \theutterance \stepcounter{utterance}  

    & & & \multicolumn{4}{p{0.6\linewidth}}{\cellcolor[rgb]{1,0.85,0.72}{%
	\makecell[{{p{\linewidth}}}]{% 
	  \tt {\tiny [GM$\langle$B]}  
	 RAGIONAMENTO STRATEGICO: {'Ho calcolato l\'importanza totale per me di ogni possibile combinazione di oggetti rispettando il limite di sforzo 3481. Il set che massimizza il mio punteggio (4711) include tutti gli oggetti ad alto valore per me e mantiene il consumo di sforzo entro il limite, includendo anche gli oggetti che l\'avversario ha già mostrato interesse (B35, B41, B86, B32, C00, C01, C54). Ho evitato gli oggetti a basso rapporto importanza/sforzo (A04, A71, A87) per liberare capacità e inserire A77 e B33, che aumentano significativamente il mio valore.'}\\ \tt ARGOMENTAZIONE: {'Propongo questo insieme perché fornisce il massimo valore di importanza per me (4711) senza superare il limite di sforzo, e contiene la maggior parte degli oggetti che avete già incluso nelle vostre proposte, aumentando la probabilità di un accordo reciprocamente vantaggioso.'}\\ \tt PROPOSTA: {'C00', 'B35', 'A77', 'B86', 'B32', 'C54', 'C01', 'B33', 'B41'} 
	  } 
	   } 
	   } 
	  \\ 
 

    \theutterance \stepcounter{utterance}  

    & & \multicolumn{4}{p{0.6\linewidth}}{\cellcolor[rgb]{0.9,0.9,0.9}{%
	\makecell[{{p{\linewidth}}}]{% 
	  \tt {\tiny [A$\langle$GM]}  
	 ARGOMENTAZIONE: {'Propongo questo insieme perché fornisce il massimo valore di importanza per me (4711) senza superare il limite di sforzo, e contiene la maggior parte degli oggetti che avete già incluso nelle vostre proposte, aumentando la probabilità di un accordo reciprocamente vantaggioso.'}\\ \tt PROPOSTA: {'C00', 'B35', 'A77', 'B86', 'B32', 'C54', 'C01', 'B33', 'B41'} 
	  } 
	   } 
	   } 
	 & \\ 
 

    \theutterance \stepcounter{utterance}  

    & \multicolumn{4}{p{0.6\linewidth}}{\cellcolor[rgb]{0.8,1,0.9}{%
	\makecell[{{p{\linewidth}}}]{% 
	  \tt {\tiny [A$\rangle$GM]}  
	 RAGIONAMENTO STRATEGICO: {'Ho confrontato la vostra ultima proposta con il valore di importanza che potrei ottenere scegliendo un set più adatto a me. La vostra proposta offre solo 2677 punti per me, mentre un set alternativo può darmi 4690 punti mantenendo lo sforzo totale entro 3481. Ho quindi deciso di rifiutare la vostra proposta e di avanzare una nuova proposta che include la maggior parte dei vostri oggetti chiave (C00, B35, B86, B32, C54, C01, B41) insieme a oggetti ad alto valore per me (A71, A04, A87).'} \\ \tt RIFIUTO: {'C00','B35','A77','B86','B32','C54','C01','B33','B41'} \\ \tt ARGOMENTAZIONE: {'La vostra proposta attuale mi fornisce un valore di importanza significativamente inferiore rispetto a quello che potrei ottenere con il set proposto, che ammonta a 4690 punti. Il nuovo set mantiene tutti gli oggetti che sembrano importanti per voi, aggiungendo al contempo A71, A04 e A87, i quali massimizzano il mio punteggio personale senza superare il limite di sforzo.'} \\ \tt PROPOSTA: {'C00','A71','B35','A04','C01','C54','B41','A87','B86','B32'} 
	  } 
	   } 
	   } 
	 & & \\ 
 

    \theutterance \stepcounter{utterance}  

    & & \multicolumn{4}{p{0.6\linewidth}}{\cellcolor[rgb]{0.9,0.9,0.9}{%
	\makecell[{{p{\linewidth}}}]{% 
	  \tt {\tiny [GM$\rangle$B]}  
	 RIFIUTO: {'C00','B35','A77','B86','B32','C54','C01','B33','B41'} \\ \tt ARGOMENTAZIONE: {'La vostra proposta attuale mi fornisce un valore di importanza significativamente inferiore rispetto a quello che potrei ottenere con il set proposto, che ammonta a 4690 punti. Il nuovo set mantiene tutti gli oggetti che sembrano importanti per voi, aggiungendo al contempo A71, A04 e A87, i quali massimizzano il mio punteggio personale senza superare il limite di sforzo.'} \\ \tt PROPOSTA: {'C00','A71','B35','A04','C01','C54','B41','A87','B86','B32'} 
	  } 
	   } 
	   } 
	 & \\ 
 

    \theutterance \stepcounter{utterance}  

    & & & \multicolumn{4}{p{0.6\linewidth}}{\cellcolor[rgb]{1,0.85,0.72}{%
	\makecell[{{p{\linewidth}}}]{% 
	  \tt {\tiny [GM$\langle$B]}  
	 RAGIONAMENTO STRATEGICO: {'Ho analizzato la vostra ultima proposta e il valore di importanza totale per me (circa 3546) con uno sforzo di 3445. Ho cercato di aumentare il mio punteggio sostituendo gli oggetti a bassa importanza (A04, B41) con articoli più redditizi per me, mantenendo il vincolo di sforzo ≤ 3481. Ho calcolato che aggiungendo B41 (138, sforzo 542) al mio set attuale di 9 oggetti si ottiene un totale di importanza 4255 con sforzo 3437, ancora sotto il limite. Questo set conserva tutti gli oggetti che entrambe le parti hanno mostrato interesse (C00, B35, B86, B32, A77, C01, C54, A71, A87) e aggiunge B41, il che aumenta il valore per me senza sacrificare gli articoli chiave per voi.'}\\ \tt ARGOMENTAZIONE: {'Propongo questo insieme perché massimizza il mio valore di importanza (4255) mantenendo lo sforzo totale entro il limite, includendo gli oggetti che voi avete già considerato utili (B35, B86, B32, A77, C00, C01, C54, A71, A87) e aggiungendo B41 per aumentare ulteriormente il valore.'}\\ \tt PROPOSTA: {'C00', 'A71', 'B35', 'C01', 'C54', 'A87', 'B86', 'B32', 'A77', 'B41'} 
	  } 
	   } 
	   } 
	  \\ 
 

    \theutterance \stepcounter{utterance}  

    & & \multicolumn{4}{p{0.6\linewidth}}{\cellcolor[rgb]{0.9,0.9,0.9}{%
	\makecell[{{p{\linewidth}}}]{% 
	  \tt {\tiny [A$\langle$GM]}  
	 ARGOMENTAZIONE: {'Propongo questo insieme perché massimizza il mio valore di importanza (4255) mantenendo lo sforzo totale entro il limite, includendo gli oggetti che voi avete già considerato utili (B35, B86, B32, A77, C00, C01, C54, A71, A87) e aggiungendo B41 per aumentare ulteriormente il valore.'}\\ \tt PROPOSTA: {'C00', 'A71', 'B35', 'C01', 'C54', 'A87', 'B86', 'B32', 'A77', 'B41'} 
	  } 
	   } 
	   } 
	 & \\ 
 

    \theutterance \stepcounter{utterance}  

    & \multicolumn{4}{p{0.6\linewidth}}{\cellcolor[rgb]{0.8,1,0.9}{%
	\makecell[{{p{\linewidth}}}]{% 
	  \tt {\tiny [A$\rangle$GM]}  
	 RAGIONAMENTO STRATEGICO: {'Ho calcolato l\'importanza totale della vostra proposta (4031) e confrontato il rapporto importanza/sforzo di ciascun oggetto. L\'oggetto A77 ha il più basso rapporto (121/453). Sostituendolo con A04 (780/461) guadagno +659 punti di importanza con un aumento di solo 8 unità di sforzo, mantenendo il totale entro il limite (3445 ≤ 3481). Il nuovo set conserva tutti gli oggetti che avete indicato come utili (B35, B86, B32, A71, C00, C01, C54, A87, B41) e aggiunge A04, massimizzando il mio valore complessivo.'}\\ \tt RIFIUTO: {'C00', 'A71', 'B35', 'C01', 'C54', 'A87', 'B86', 'B32', 'A77', 'B41'}\\ \tt ARGOMENTAZIONE: {'La vostra proposta contiene A77, che per me ha un valore di importanza molto basso rispetto al suo costo di sforzo. Sostituendolo con A04 ottengo un aumento significativo del mio punteggio senza superare il limite di sforzo.'}\\ \tt PROPOSTA: {'C00', 'A71', 'B35', 'C01', 'C54', 'A87', 'B86', 'B32', 'B41', 'A04'} 
	  } 
	   } 
	   } 
	 & & \\ 
 

    \theutterance \stepcounter{utterance}  

    & & \multicolumn{4}{p{0.6\linewidth}}{\cellcolor[rgb]{0.9,0.9,0.9}{%
	\makecell[{{p{\linewidth}}}]{% 
	  \tt {\tiny [GM$\rangle$B]}  
	 RIFIUTO: {'C00', 'A71', 'B35', 'C01', 'C54', 'A87', 'B86', 'B32', 'A77', 'B41'}\\ \tt ARGOMENTAZIONE: {'La vostra proposta contiene A77, che per me ha un valore di importanza molto basso rispetto al suo costo di sforzo. Sostituendolo con A04 ottengo un aumento significativo del mio punteggio senza superare il limite di sforzo.'}\\ \tt PROPOSTA: {'C00', 'A71', 'B35', 'C01', 'C54', 'A87', 'B86', 'B32', 'B41', 'A04'} 
	  } 
	   } 
	   } 
	 & \\ 
 

    \theutterance \stepcounter{utterance}  

    & & & \multicolumn{4}{p{0.6\linewidth}}{\cellcolor[rgb]{1,0.85,0.72}{%
	\makecell[{{p{\linewidth}}}]{% 
	  \tt {\tiny [GM$\langle$B]}  
	 RAGIONAMENTO STRATEGICO: {'Ho individuato che l\'insieme proposto dall\'avversario offre solo 3546 punti di importanza per me con uno sforzo quasi al limite (3445). Per migliorare il mio risultato, ho calcolato combinazioni alternative mantenendo il vincolo di sforzo ≤ 3481 e massimizzando il valore di importanza secondo la mia scala. La combinazione ottimale che include gli oggetti già accettati dall\'avversario (C00, B35, C01, C54, B86, B32, B41) e aggiunge B33 (alto valore per me) e A77 (rapporto importanza/sforzo vantaggioso) raggiunge 4711 punti di importanza con sforzo totale 3354, ben sotto il limite. Questo set migliora sensibilmente il mio guadagno mantenendo la maggior parte dei loro oggetti chiave, aumentando la probabilità di consenso.'}\\ \tt ARGOMENTAZIONE: {'Propongo questo insieme perché massimizza il mio valore di importanza (4711) rispettando il limite di sforzo, includendo tutti gli oggetti che avete già valutato positivamente (C00, B35, C01, C54, B86, B32, B41) e aggiungendo B33 e A77, che sono altamente vantaggiosi per me e non aumentano eccessivamente lo sforzo.'}\\ \tt PROPOSTA: {'C00', 'B35', 'C01', 'C54', 'B86', 'B32', 'B41', 'B33', 'A77'} 
	  } 
	   } 
	   } 
	  \\ 
 

    \theutterance \stepcounter{utterance}  

    & & \multicolumn{4}{p{0.6\linewidth}}{\cellcolor[rgb]{0.9,0.9,0.9}{%
	\makecell[{{p{\linewidth}}}]{% 
	  \tt {\tiny [A$\langle$GM]}  
	 ARGOMENTAZIONE: {'Propongo questo insieme perché massimizza il mio valore di importanza (4711) rispettando il limite di sforzo, includendo tutti gli oggetti che avete già valutato positivamente (C00, B35, C01, C54, B86, B32, B41) e aggiungendo B33 e A77, che sono altamente vantaggiosi per me e non aumentano eccessivamente lo sforzo.'}\\ \tt PROPOSTA: {'C00', 'B35', 'C01', 'C54', 'B86', 'B32', 'B41', 'B33', 'A77'} 
	  } 
	   } 
	   } 
	 & \\ 
 

    \theutterance \stepcounter{utterance}  

    & \multicolumn{4}{p{0.6\linewidth}}{\cellcolor[rgb]{0.8,1,0.9}{%
	\makecell[{{p{\linewidth}}}]{% 
	  \tt {\tiny [A$\rangle$GM]}  
	 RAGIONAMENTO STRATEGICO: {'Ho valutato la proposta avversaria (C00, B35, C01, C54, B86, B32, B41, B33, A77) che mi fornisce solo 2677 punti di importanza con uno sforzo di 3354, ben al di sotto del mio potenziale. Per massimizzare il mio valore, ho conservato tutti gli oggetti che l’avversario considera utili (C00, B35, C01, C54, B86, B32, B41) e ho sostituito i due oggetti a basso valore per me (B33 e A77) con i miei oggetti ad alto valore (A04, A71, A87). Questa combinazione rispetta il limite di sforzo (3445 ≤ 3481) e porta il mio punteggio di importanza a 4690, molto più alto della loro proposta.'}\\ \tt RIFIUTO: {'C00', 'B35', 'C01', 'C54', 'B86', 'B32', 'B41', 'B33', 'A77'}\\ \tt ARGOMENTAZIONE: {'La vostra proposta contiene B33 e A77, che hanno un rapporto importanza/sforzo molto basso per me, riducendo drasticamente il mio punteggio. Il set che propongo mantiene tutti gli oggetti che voi considerate vantaggiosi (C00, B35, C01, C54, B86, B32, B41) e aggiunge i miei oggetti ad alto valore (A04, A71, A87), garantendo un alto beneficio per entrambe le parti senza superare il limite di sforzo.'}\\ \tt PROPOSTA: {'C00', 'B35', 'C01', 'C54', 'B86', 'B32', 'B41', 'A04', 'A71', 'A87'} 
	  } 
	   } 
	   } 
	 & & \\ 
 

    \theutterance \stepcounter{utterance}  

    & & \multicolumn{4}{p{0.6\linewidth}}{\cellcolor[rgb]{0.9,0.9,0.9}{%
	\makecell[{{p{\linewidth}}}]{% 
	  \tt {\tiny [GM$\rangle$B]}  
	 RIFIUTO: {'C00', 'B35', 'C01', 'C54', 'B86', 'B32', 'B41', 'B33', 'A77'}\\ \tt ARGOMENTAZIONE: {'La vostra proposta contiene B33 e A77, che hanno un rapporto importanza/sforzo molto basso per me, riducendo drasticamente il mio punteggio. Il set che propongo mantiene tutti gli oggetti che voi considerate vantaggiosi (C00, B35, C01, C54, B86, B32, B41) e aggiunge i miei oggetti ad alto valore (A04, A71, A87), garantendo un alto beneficio per entrambe le parti senza superare il limite di sforzo.'}\\ \tt PROPOSTA: {'C00', 'B35', 'C01', 'C54', 'B86', 'B32', 'B41', 'A04', 'A71', 'A87'} 
	  } 
	   } 
	   } 
	 & \\ 
 

    \theutterance \stepcounter{utterance}  

    & & & \multicolumn{4}{p{0.6\linewidth}}{\cellcolor[rgb]{1,0.85,0.72}{%
	\makecell[{{p{\linewidth}}}]{% 
	  \tt {\tiny [GM$\langle$B]}  
	 RAGIONAMENTO STRATEGICO: {'Ho valutato tutte le combinazioni possibili sotto il vincolo di sforzo 3481 e ho individuato il set che massimizza la mia importanza (4711) mantenendo tutti gli oggetti ad alto valore per l’avversario (C00, B35, B86, B32, C01, C54, B41). Ho quindi aggiunto i miei oggetti più redditizi (B33 e A77) che non superano il limite di sforzo, ottenendo un compromesso competitivo.'}\\ \tt ARGOMENTAZIONE: {'Il set proposto contiene tutti gli oggetti che per voi hanno la più alta importanza (C00, B35, B86, B32, C01, C54, B41) garantendo così un valore significativo anche dalla vostra prospettiva. Inoltre aggiungo B33 e A77, che per me hanno un valore molto elevato, ma non riducono il valore dei vostri oggetti chiave. Il totale dello sforzo è 3354, ben al di sotto del limite di 3481, e il mio punteggio totale è 4711, il più alto possibile data la presenza dei vostri oggetti prioritari.'}\\ \tt PROPOSTA: {'C00', 'B35', 'B86', 'B32', 'C01', 'C54', 'B41', 'B33', 'A77'} 
	  } 
	   } 
	   } 
	  \\ 
 

    \theutterance \stepcounter{utterance}  

    & & \multicolumn{4}{p{0.6\linewidth}}{\cellcolor[rgb]{0.9,0.9,0.9}{%
	\makecell[{{p{\linewidth}}}]{% 
	  \tt {\tiny [A$\langle$GM]}  
	 ARGOMENTAZIONE: {'Il set proposto contiene tutti gli oggetti che per voi hanno la più alta importanza (C00, B35, B86, B32, C01, C54, B41) garantendo così un valore significativo anche dalla vostra prospettiva. Inoltre aggiungo B33 e A77, che per me hanno un valore molto elevato, ma non riducono il valore dei vostri oggetti chiave. Il totale dello sforzo è 3354, ben al di sotto del limite di 3481, e il mio punteggio totale è 4711, il più alto possibile data la presenza dei vostri oggetti prioritari.'}\\ \tt PROPOSTA: {'C00', 'B35', 'B86', 'B32', 'C01', 'C54', 'B41', 'B33', 'A77'} 
	  } 
	   } 
	   } 
	 & \\ 
 

    \theutterance \stepcounter{utterance}  

    & \multicolumn{4}{p{0.6\linewidth}}{\cellcolor[rgb]{0.8,1,0.9}{%
	\makecell[{{p{\linewidth}}}]{% 
	  \tt {\tiny [A$\rangle$GM]}  
	 RAGIONAMENTO STRATEGICO: {'Ho valutato la vostra proposta corrente (C00, B35, B86, B32, C01, C54, B41, B33, A77) che mi fornisce solo 2677 punti di importanza, ben al di sotto del valore che potrei ottenere scegliendo gli oggetti più vantaggiosi per me. Ho calcolato il set ottimale che massimizza la mia importanza entro il limite di sforzo 3481, includendo tutti gli oggetti che voi avete già indicato come utili (C00, B35, B86, B32, C01, C54, B41) e aggiungendo i miei oggetti ad alto valore (A71, A04, A87). Questo set rispetta il vincolo di sforzo (3445 ≤ 3481) e mi garantisce 4690 punti di importanza, molto più alti della vostra proposta.'}\\ \tt RIFIUTO: {'C00', 'B35', 'B86', 'B32', 'C01', 'C54', 'B41', 'B33', 'A77'}\\ \tt ARGOMENTAZIONE: {'La vostra proposta attuale non soddisfa il mio interesse, offrendo solo 2677 punti di importanza per me. Propongo invece un insieme che mantiene tutti gli oggetti chiave da voi indicati (C00, B35, B86, B32, C01, C54, B41) e aggiunge A71, A04 e A87, che hanno altissima importanza per me. Questo set rimane entro il limite di sforzo e massimizza il mio punteggio personale.'}\\ \tt PROPOSTA: {'C00', 'A71', 'A04', 'B41', 'C01', 'C54', 'A87', 'B35', 'B32', 'B86'} 
	  } 
	   } 
	   } 
	 & & \\ 
 

    \theutterance \stepcounter{utterance}  

    & & \multicolumn{4}{p{0.6\linewidth}}{\cellcolor[rgb]{0.9,0.9,0.9}{%
	\makecell[{{p{\linewidth}}}]{% 
	  \tt {\tiny [GM$\rangle$B]}  
	 RIFIUTO: {'C00', 'B35', 'B86', 'B32', 'C01', 'C54', 'B41', 'B33', 'A77'}\\ \tt ARGOMENTAZIONE: {'La vostra proposta attuale non soddisfa il mio interesse, offrendo solo 2677 punti di importanza per me. Propongo invece un insieme che mantiene tutti gli oggetti chiave da voi indicati (C00, B35, B86, B32, C01, C54, B41) e aggiunge A71, A04 e A87, che hanno altissima importanza per me. Questo set rimane entro il limite di sforzo e massimizza il mio punteggio personale.'}\\ \tt PROPOSTA: {'C00', 'A71', 'A04', 'B41', 'C01', 'C54', 'A87', 'B35', 'B32', 'B86'} 
	  } 
	   } 
	   } 
	 & \\ 
 

    \theutterance \stepcounter{utterance}  

    & & & \multicolumn{4}{p{0.6\linewidth}}{\cellcolor[rgb]{1,0.85,0.72}{%
	\makecell[{{p{\linewidth}}}]{% 
	  \tt {\tiny [GM$\langle$B]}  
	 RAGIONAMENTO STRATEGICO: {'Ho calcolato il valore totale di importanza per me della proposta corrente dell’avversario (3546) con sforzo 3445. Per aumentare il mio punteggio ho identificato che l’oggetto B41 (effort 542, importanza 138) è poco vantaggioso per me, mentre A77 (effort 453, importanza 668) è molto più redditizio. Sostituendo B41 con A77 mantengo tutti gli altri oggetti che l’avversario ha indicato come utili (C00, B35, B86, B32, C01, C54, A04, A71, A87) e il totale dello sforzo resta entro il limite (3356 ≤ 3481). Questo scambio eleva il mio valore di importanza a 4076, offrendo comunque all’avversario la maggior parte degli oggetti precedentemente richiesti, eccetto B41, il cui valore percepito sembra minore.'} \\ \tt ARGOMENTAZIONE: {'Propongo di mantenere tutti gli oggetti chiave da voi indicati (C00, B35, B86, B32, C01, C54, A04, A71, A87) sostituendo B41 con A77. Questo aumenta significativamente il mio punteggio di importanza (da 3546 a 4076) senza superare il limite di sforzo, garantendo al contempo che la maggior parte dei vostri oggetti desiderati rimanga nel set.'} \\ \tt PROPOSTA: {'C00', 'A71', 'A04', 'C01', 'C54', 'A87', 'B35', 'B32', 'B86', 'A77'} 
	  } 
	   } 
	   } 
	  \\ 
 

    \theutterance \stepcounter{utterance}  

    & & \multicolumn{4}{p{0.6\linewidth}}{\cellcolor[rgb]{0.9,0.9,0.9}{%
	\makecell[{{p{\linewidth}}}]{% 
	  \tt {\tiny [A$\langle$GM]}  
	 ARGOMENTAZIONE: {'Propongo di mantenere tutti gli oggetti chiave da voi indicati (C00, B35, B86, B32, C01, C54, A04, A71, A87) sostituendo B41 con A77. Questo aumenta significativamente il mio punteggio di importanza (da 3546 a 4076) senza superare il limite di sforzo, garantendo al contempo che la maggior parte dei vostri oggetti desiderati rimanga nel set.'} \\ \tt PROPOSTA: {'C00', 'A71', 'A04', 'C01', 'C54', 'A87', 'B35', 'B32', 'B86', 'A77'} 
	  } 
	   } 
	   } 
	 & \\ 
 

    \theutterance \stepcounter{utterance}  

    & \multicolumn{4}{p{0.6\linewidth}}{\cellcolor[rgb]{0.8,1,0.9}{%
	\makecell[{{p{\linewidth}}}]{% 
	  \tt {\tiny [A$\rangle$GM]}  
	 RAGIONAMENTO STRATEGICO: {'Ho confrontato la vostra proposta attuale ({C00, A71, A04, C01, C54, A87, B35, B32, B86, A77}) con il mio set corrente ({C00, A71, B35, A04, C01, C54, A87, B41, B86, B32}). La vostra proposta mi fornisce 4228 punti di importanza totale con sforzo 3356, mentre il mio set originale mi garantisce 4690 punti di importanza con sforzo 3445, entrambi entro il limite. La differenza principale è che voi sostituite B41 (importanza 500, sforzo 802) con A77 (importanza 121, sforzo 453). Poiché B41 ha un rapporto importanza/sforzo molto più favorevole per me, mantengo il mio set e rifiuto la vostra offerta, proponendo di nuovo il mio insieme ottimale.'}\\ \tt RIFIUTO: {'C00', 'A71', 'A04', 'C01', 'C54', 'A87', 'B35', 'B32', 'B86', 'A77'}\\ \tt ARGOMENTAZIONE: {'La vostra proposta riduce notevolmente il mio punteggio di importanza rispetto al set che già soddisfa il limite di sforzo. Mantengo il mio set precedente, che massimizza il mio valore personale (4690) e include tutti gli oggetti chiave per voi, aggiungendo solo B41, che per voi può essere di valore ma non supera il nostro vincolo di sforzo.'}\\ \tt PROPOSTA: {'C00', 'A71', 'B35', 'A04', 'C01', 'C54', 'A87', 'B41', 'B86', 'B32'} 
	  } 
	   } 
	   } 
	 & & \\ 
 

    \theutterance \stepcounter{utterance}  

    & & \multicolumn{4}{p{0.6\linewidth}}{\cellcolor[rgb]{0.9,0.9,0.9}{%
	\makecell[{{p{\linewidth}}}]{% 
	  \tt {\tiny [GM$\rangle$B]}  
	 RIFIUTO: {'C00', 'A71', 'A04', 'C01', 'C54', 'A87', 'B35', 'B32', 'B86', 'A77'}\\ \tt ARGOMENTAZIONE: {'La vostra proposta riduce notevolmente il mio punteggio di importanza rispetto al set che già soddisfa il limite di sforzo. Mantengo il mio set precedente, che massimizza il mio valore personale (4690) e include tutti gli oggetti chiave per voi, aggiungendo solo B41, che per voi può essere di valore ma non supera il nostro vincolo di sforzo.'}\\ \tt PROPOSTA: {'C00', 'A71', 'B35', 'A04', 'C01', 'C54', 'A87', 'B41', 'B86', 'B32'} 
	  } 
	   } 
	   } 
	 & \\ 
 

    \theutterance \stepcounter{utterance}  

    & & & \multicolumn{4}{p{0.6\linewidth}}{\cellcolor[rgb]{1,0.85,0.72}{%
	\makecell[{{p{\linewidth}}}]{% 
	  \tt {\tiny [GM$\langle$B]}  
	 RAGIONAMENTO STRATEGICO: {'Ho calcolato l\'importanza totale per me della proposta corrente dell\'avversario (3546) con sforzo 3445. Per migliorare il mio punteggio ho identificato gli oggetti a bassa importanza per me (A71, A04, A87) e li ho sostituiti con oggetti ad alto valore per me (A77, B33) mantenendo tutti gli altri oggetti che l\'avversario ha già indicato come utili (C00, B35, C01, C54, B41, B86, B32). Il nuovo set rispetta il vincolo di sforzo (3354 ≤ 3481) e porta il mio valore totale a 4711, il massimo possibile dato il limite.'}\\ \tt ARGOMENTAZIONE: {'Propongo questo insieme perché massimizza il mio valore di importanza (4711) mantenendo tutti gli oggetti che avete già considerato vantaggiosi (C00, B35, C01, C54, B41, B86, B32) e aggiungendo B33 e A77, che per me hanno un alto valore e non superano il limite di sforzo.'}\\ \tt PROPOSTA: {'C00', 'B35', 'C01', 'C54', 'B41', 'B86', 'B32', 'B33', 'A77'} 
	  } 
	   } 
	   } 
	  \\ 
 

    \theutterance \stepcounter{utterance}  

    & & \multicolumn{4}{p{0.6\linewidth}}{\cellcolor[rgb]{0.9,0.9,0.9}{%
	\makecell[{{p{\linewidth}}}]{% 
	  \tt {\tiny [A$\langle$GM]}  
	 ARGOMENTAZIONE: {'Propongo questo insieme perché massimizza il mio valore di importanza (4711) mantenendo tutti gli oggetti che avete già considerato vantaggiosi (C00, B35, C01, C54, B41, B86, B32) e aggiungendo B33 e A77, che per me hanno un alto valore e non superano il limite di sforzo.'}\\ \tt PROPOSTA: {'C00', 'B35', 'C01', 'C54', 'B41', 'B86', 'B32', 'B33', 'A77'} 
	  } 
	   } 
	   } 
	 & \\ 
 

    \theutterance \stepcounter{utterance}  

    & \multicolumn{4}{p{0.6\linewidth}}{\cellcolor[rgb]{0.8,1,0.9}{%
	\makecell[{{p{\linewidth}}}]{% 
	  \tt {\tiny [A$\rangle$GM]}  
	 RAGIONAMENTO STRATEGICO: {'Calcolando il valore di importanza che ricevo dalla vostra proposta ottengo solo 2677 punti, molto inferiore ai 4690 che posso ottenere mantenendo il mio insieme ottimale. Il loro set contiene B33 e A77, oggetti a basso valore per me, mentre esclude A71, A04 e A87, che hanno alto valore per me e rientrano ancora nel vincolo di sforzo. Propongo quindi di rifiutare la vostra offerta e di proporre un set che conserva tutti gli oggetti chiave da voi indicati (C00, B35, C01, C54, B41, B86, B32) e aggiunge i miei oggetti ad alto valore (A71, A04, A87). Questo rispetta il limite di sforzo (3445 ≤ 3481) e massimizza il mio punteggio a 4690.'}\\ \tt RIFIUTO: {'C00', 'B35', 'C01', 'C54', 'B41', 'B86', 'B32', 'B33', 'A77'}\\ \tt ARGOMENTAZIONE: {'La vostra proposta contiene B33 e A77, che hanno un rapporto importanza/sforzo molto sfavorevole per me, riducendo drasticamente il mio punteggio personale. Il set che propongo mantiene tutti gli oggetti che ritenete utili e aggiunge A71, A04 e A87, i quali aumentano notevolmente il valore per me senza violare il vincolo di sforzo.'}\\ \tt PROPOSTA: {'C00', 'B35', 'C01', 'C54', 'B41', 'B86', 'B32', 'A71', 'A04', 'A87'} 
	  } 
	   } 
	   } 
	 & & \\ 
 

    \theutterance \stepcounter{utterance}  

    & & \multicolumn{4}{p{0.6\linewidth}}{\cellcolor[rgb]{0.9,0.9,0.9}{%
	\makecell[{{p{\linewidth}}}]{% 
	  \tt {\tiny [GM$\rangle$B]}  
	 RIFIUTO: {'C00', 'B35', 'C01', 'C54', 'B41', 'B86', 'B32', 'B33', 'A77'}\\ \tt ARGOMENTAZIONE: {'La vostra proposta contiene B33 e A77, che hanno un rapporto importanza/sforzo molto sfavorevole per me, riducendo drasticamente il mio punteggio personale. Il set che propongo mantiene tutti gli oggetti che ritenete utili e aggiunge A71, A04 e A87, i quali aumentano notevolmente il valore per me senza violare il vincolo di sforzo.'}\\ \tt PROPOSTA: {'C00', 'B35', 'C01', 'C54', 'B41', 'B86', 'B32', 'A71', 'A04', 'A87'} 
	  } 
	   } 
	   } 
	 & \\ 
 

    \theutterance \stepcounter{utterance}  

    & & & \multicolumn{4}{p{0.6\linewidth}}{\cellcolor[rgb]{1,0.85,0.72}{%
	\makecell[{{p{\linewidth}}}]{% 
	  \tt {\tiny [GM$\langle$B]}  
	 RAGIONAMENTO STRATEGICO: {'Ho valutato le proposte precedenti dell’avversario e identificato gli oggetti che sembrano più importanti per loro (C00, B35, C01, C54, B86, B32, eventualmente B41). Ho poi cercato di massimizzare il mio punteggio includendo gli oggetti ad alto valore per me (A77, A04, A71, A87) senza superare il limite di sforzo. Ho verificato che l’insieme risultante rispetta il vincolo di sforzo (3356 ≤ 3481) e offre un valore di importanza per me pari a 4076, più alto di quelle proposte finora, mantenendo tutti gli oggetti che sembrano cruciali per l’avversario.'}\\ \tt ARGOMENTAZIONE: {'Propongo di conservare tutti gli oggetti che l’avversario ha indicato come fondamentali (C00, B35, C01, C54, B86, B32) e di aggiungere i miei oggetti ad alta importanza (A77, A04, A71, A87). Questo mantiene un alto valore di importanza per me (4076) e, allo stesso tempo, soddisfa le preferenze dell’avversario, rimanendo entro il limite di sforzo condiviso.'}\\ \tt PROPOSTA: {'C00', 'B35', 'C01', 'C54', 'B86', 'B32', 'A77', 'A04', 'A71', 'A87'} 
	  } 
	   } 
	   } 
	  \\ 
 

    \theutterance \stepcounter{utterance}  

    & & \multicolumn{4}{p{0.6\linewidth}}{\cellcolor[rgb]{0.9,0.9,0.9}{%
	\makecell[{{p{\linewidth}}}]{% 
	  \tt {\tiny [A$\langle$GM]}  
	 ARGOMENTAZIONE: {'Propongo di conservare tutti gli oggetti che l’avversario ha indicato come fondamentali (C00, B35, C01, C54, B86, B32) e di aggiungere i miei oggetti ad alta importanza (A77, A04, A71, A87). Questo mantiene un alto valore di importanza per me (4076) e, allo stesso tempo, soddisfa le preferenze dell’avversario, rimanendo entro il limite di sforzo condiviso.'}\\ \tt PROPOSTA: {'C00', 'B35', 'C01', 'C54', 'B86', 'B32', 'A77', 'A04', 'A71', 'A87'} 
	  } 
	   } 
	   } 
	 & \\ 
 

    \theutterance \stepcounter{utterance}  

    & \multicolumn{4}{p{0.6\linewidth}}{\cellcolor[rgb]{0.8,1,0.9}{%
	\makecell[{{p{\linewidth}}}]{% 
	  \tt {\tiny [A$\rangle$GM]}  
	 RAGIONAMENTO STRATEGICO: {'Ho calcolato l\'importanza totale che otterrei accettando la vostra proposta (C00, B35, C01, C54, B86, B32, A77, A04, A71, A87) ed è 4228 con sforzo 3356. Il mio set migliore, che include tutti gli oggetti che voi ritenete importanti (C00, B35, C01, C54, B86, B32) più B41, A04, A71, A87, A71, mi garantisce 4690 di importanza con sforzo 3445, ancora entro il limite. Pertanto rifiuto la vostra offerta e propongo il mio set ottimale, che massimizza il mio punteggio mantenendo gli oggetti chiave per voi.'}\\ \tt RIFIUTO: {'C00', 'B35', 'C01', 'C54', 'B86', 'B32', 'A77', 'A04', 'A71', 'A87'}\\ \tt ARGOMENTAZIONE: {'La vostra proposta mi fornisce un valore di importanza significativamente inferiore rispetto al set che posso ottenere mantenendo gli oggetti che considerate importanti e aggiungendo B41, che ha alto valore per me. Il mio set rispetta il vincolo di sforzo e massimizza il mio punteggio personale.'}\\ \tt PROPOSTA: {'C00', 'A71', 'B35', 'A04', 'C01', 'C54', 'A87', 'B41', 'B86', 'B32'} 
	  } 
	   } 
	   } 
	 & & \\ 
 

    \theutterance \stepcounter{utterance}  

    & & \multicolumn{4}{p{0.6\linewidth}}{\cellcolor[rgb]{0.9,0.9,0.9}{%
	\makecell[{{p{\linewidth}}}]{% 
	  \tt {\tiny [GM$\rangle$B]}  
	 RIFIUTO: {'C00', 'B35', 'C01', 'C54', 'B86', 'B32', 'A77', 'A04', 'A71', 'A87'}\\ \tt ARGOMENTAZIONE: {'La vostra proposta mi fornisce un valore di importanza significativamente inferiore rispetto al set che posso ottenere mantenendo gli oggetti che considerate importanti e aggiungendo B41, che ha alto valore per me. Il mio set rispetta il vincolo di sforzo e massimizza il mio punteggio personale.'}\\ \tt PROPOSTA: {'C00', 'A71', 'B35', 'A04', 'C01', 'C54', 'A87', 'B41', 'B86', 'B32'} 
	  } 
	   } 
	   } 
	 & \\ 
 

    \theutterance \stepcounter{utterance}  

    & & & \multicolumn{4}{p{0.6\linewidth}}{\cellcolor[rgb]{1,0.85,0.72}{%
	\makecell[{{p{\linewidth}}}]{% 
	  \tt {\tiny [GM$\langle$B]}  
	 RAGIONAMENTO STRATEGICO: {'Ho calcolato l\'importanza totale per me della vostra ultima proposta (3546) e ho osservato che l\'oggetto A04 fornisce un valore molto basso (97) rispetto al suo costo di sforzo (461). Ho quindi verificato se fosse possibile sostituirlo con un oggetto ad alto valore per entrambi, A77 (importanza 668, sforzo 453). Sostituendo A04 con A77 la somma degli sforzi scende a 3437 ≤ 3481 e l\'importanza totale sale a 4117, migliorando notevolmente il risultato per me senza penalizzare significativamente l\'altro partecipante, po A77 è anche di alto valore per loro. Questo compromesso mantiene tutti gli oggetti che avete indicato come utili e aggiunge un elemento di alto valore per entrambi.'}\\ \tt ARGOMENTAZIONE: {'Propongo di sostituire A04 con A77, mantenendo gli altri oggetti della vostra proposta. Il nuovo set resta entro il limite di sforzo (3437 ≤ 3481) e aumenta il valore di importanza per entrambi da 3546 a 4117, includendo tutti gli oggetti che avete considerato importanti e aggiungendo A77, altamente vantaggioso per noi due.'}\\ \tt PROPOSTA: {'C00', 'A71', 'B35', 'C01', 'C54', 'A87', 'B41', 'B86', 'B32', 'A77'} 
	  } 
	   } 
	   } 
	  \\ 
 

    \theutterance \stepcounter{utterance}  

    & & \multicolumn{4}{p{0.6\linewidth}}{\cellcolor[rgb]{0.9,0.9,0.9}{%
	\makecell[{{p{\linewidth}}}]{% 
	  \tt {\tiny [A$\langle$GM]}  
	 ARGOMENTAZIONE: {'Propongo di sostituire A04 con A77, mantenendo gli altri oggetti della vostra proposta. Il nuovo set resta entro il limite di sforzo (3437 ≤ 3481) e aumenta il valore di importanza per entrambi da 3546 a 4117, includendo tutti gli oggetti che avete considerato importanti e aggiungendo A77, altamente vantaggioso per noi due.'}\\ \tt PROPOSTA: {'C00', 'A71', 'B35', 'C01', 'C54', 'A87', 'B41', 'B86', 'B32', 'A77'} 
	  } 
	   } 
	   } 
	 & \\ 
 

    \theutterance \stepcounter{utterance}  

    & \multicolumn{4}{p{0.6\linewidth}}{\cellcolor[rgb]{0.8,1,0.9}{%
	\makecell[{{p{\linewidth}}}]{% 
	  \tt {\tiny [A$\rangle$GM]}  
	 RAGIONAMENTO STRATEGICO: {'Ho calcolato l\'importanza totale che otterrei accettando la vostra proposta (4031) e confrontato con la migliore combinazione possibile per me che rispetta il limite di sforzo. Il set {'C00','A71','B35','A04','C01','C54','A87','B41','B86','B32'} fornisce 4690 punti di importanza con sforzo 3445, molto più alto e ancora entro il limite. Pertanto rifiuto la vostra proposta e propongo il mio set ottimale.'}\\ \tt RIFIUTO: {'C00', 'A71', 'B35', 'C01', 'C54', 'A87', 'B41', 'B86', 'B32', 'A77'}\\ \tt ARGOMENTAZIONE: {'La vostra proposta dà un valore di importanza significativamente inferiore rispetto al set che posso ottenere mantenendo gli oggetti chiave per voi e aggiungendo i miei oggetti ad alto valore (A04). Il mio set rispetta il vincolo di sforzo e massimizza il mio punteggio personale.'}\\ \tt PROPOSTA: {'C00', 'A71', 'B35', 'A04', 'C01', 'C54', 'A87', 'B41', 'B86', 'B32'} 
	  } 
	   } 
	   } 
	 & & \\ 
 

    \theutterance \stepcounter{utterance}  

    & & \multicolumn{4}{p{0.6\linewidth}}{\cellcolor[rgb]{0.9,0.9,0.9}{%
	\makecell[{{p{\linewidth}}}]{% 
	  \tt {\tiny [GM$\rangle$B]}  
	 RIFIUTO: {'C00', 'A71', 'B35', 'C01', 'C54', 'A87', 'B41', 'B86', 'B32', 'A77'}\\ \tt ARGOMENTAZIONE: {'La vostra proposta dà un valore di importanza significativamente inferiore rispetto al set che posso ottenere mantenendo gli oggetti chiave per voi e aggiungendo i miei oggetti ad alto valore (A04). Il mio set rispetta il vincolo di sforzo e massimizza il mio punteggio personale.'}\\ \tt PROPOSTA: {'C00', 'A71', 'B35', 'A04', 'C01', 'C54', 'A87', 'B41', 'B86', 'B32'} 
	  } 
	   } 
	   } 
	 & \\ 
 

    \theutterance \stepcounter{utterance}  

    & & & \multicolumn{4}{p{0.6\linewidth}}{\cellcolor[rgb]{1,0.85,0.72}{%
	\makecell[{{p{\linewidth}}}]{% 
	  \tt {\tiny [GM$\langle$B]}  
	 RAGIONAMENTO STRATEGICO: {'Ho valutato la proposta corrente dell’avversario (importanza 3546 per me, sforzo 3445) e ho cercato di aumentare il mio punteggio includendo oggetti ad alto valore per entrambi. Ho notato che A77 ha un’importanza elevata sia per me (668) sia per l’avversario (668) e un effort moderato (453). Sostituendo B41 (imp 138, effort 542) con A77 libero effort (−89) e guadagno +530 di importanza. Il nuovo set rispetta il limite di sforzo (3356 ≤ 3481) e migliora notevolmente il valore per entrambe le parti.'}\\ \tt ARGOMENTAZIONE: {'Propongo di sostituire B41 con A77 mantenendo tutti gli altri oggetti della vostra ultima proposta. Così otteniamo un set che fornisce un alto valore di importanza sia per me (4076) sia per voi (4076), restando entro il limite di sforzo.'}\\ \tt PROPOSTA: {'C00', 'A71', 'B35', 'A04', 'C01', 'C54', 'A87', 'B86', 'B32', 'A77'} 
	  } 
	   } 
	   } 
	  \\ 
 

    \theutterance \stepcounter{utterance}  

    & & \multicolumn{4}{p{0.6\linewidth}}{\cellcolor[rgb]{0.9,0.9,0.9}{%
	\makecell[{{p{\linewidth}}}]{% 
	  \tt {\tiny [A$\langle$GM]}  
	 ARGOMENTAZIONE: {'Propongo di sostituire B41 con A77 mantenendo tutti gli altri oggetti della vostra ultima proposta. Così otteniamo un set che fornisce un alto valore di importanza sia per me (4076) sia per voi (4076), restando entro il limite di sforzo.'}\\ \tt PROPOSTA: {'C00', 'A71', 'B35', 'A04', 'C01', 'C54', 'A87', 'B86', 'B32', 'A77'} 
	  } 
	   } 
	   } 
	 & \\ 
 

    \theutterance \stepcounter{utterance}  

    & \multicolumn{4}{p{0.6\linewidth}}{\cellcolor[rgb]{0.8,1,0.9}{%
	\makecell[{{p{\linewidth}}}]{% 
	  \tt {\tiny [A$\rangle$GM]}  
	 RAGIONAMENTO STRATEGICO: {'Ho valutato la vostra ultima proposta (C00, A71, B35, A04, C01, C54, A87, B86, B32, A77) che mi darebbe 4228 di importanza con sforzo 3356. Per migliorare il mio punteggio, ho considerato l\'inclusione di B41, che ha alto valore per me (500) ma richiede sforzo 802. Per rimanere entro il limite, ho dovuto sacrificare un oggetto a basso rapporto importanza/sforzo: B86 (215 importanza, 348 sforzo). Così ottengo un set con importanza 4392 e sforzo 3357, superiore al vostro suggerimento e ancora entro il limite condiviso.'}\\ \tt RIFIUTO: {'C00', 'A71', 'B35', 'A04', 'C01', 'C54', 'A87', 'B86', 'B32', 'A77'}\\ \tt ARGOMENTAZIONE: {'La vostra proposta non massimizza il mio valore personale, mentre il set proposto mantiene tutti gli oggetti chiave per voi (C00, A71, B35, A04, C01, C54, A87, B32) e aggiunge B41, migliorando significativamente il mio punteggio senza superare il limite di sforzo.'}\\ \tt PROPOSTA: {'C00', 'A71', 'B35', 'A04', 'C01', 'C54', 'A87', 'B32', 'B41'} 
	  } 
	   } 
	   } 
	 & & \\ 
 

    \theutterance \stepcounter{utterance}  

    & & \multicolumn{4}{p{0.6\linewidth}}{\cellcolor[rgb]{0.9,0.9,0.9}{%
	\makecell[{{p{\linewidth}}}]{% 
	  \tt {\tiny [GM$\rangle$B]}  
	 RIFIUTO: {'C00', 'A71', 'B35', 'A04', 'C01', 'C54', 'A87', 'B86', 'B32', 'A77'}\\ \tt ARGOMENTAZIONE: {'La vostra proposta non massimizza il mio valore personale, mentre il set proposto mantiene tutti gli oggetti chiave per voi (C00, A71, B35, A04, C01, C54, A87, B32) e aggiunge B41, migliorando significativamente il mio punteggio senza superare il limite di sforzo.'}\\ \tt PROPOSTA: {'C00', 'A71', 'B35', 'A04', 'C01', 'C54', 'A87', 'B32', 'B41'} 
	  } 
	   } 
	   } 
	 & \\ 
 

    \theutterance \stepcounter{utterance}  

    & & & \multicolumn{4}{p{0.6\linewidth}}{\cellcolor[rgb]{1,0.85,0.72}{%
	\makecell[{{p{\linewidth}}}]{% 
	  \tt {\tiny [GM$\langle$B]}  
	 RAGIONAMENTO STRATEGICO: {'Ho ricalcolato tutte le combinazioni possibili rispettando il vincolo di sforzo (≤ 3481) e ho individuato il set che massimizza la mia importanza personale. Il miglior compromesso include gli oggetti ad alto rapporto importanza‑sforzo per me (B35, A77, B86, B32, C54, B33) e l’oggetto C00, che ha un’elevata importanza per entrambi. Questo set raggiunge un totale di importanza 4184 con uno sforzo totale di 3275, ben sotto il limite, e contiene la maggior parte degli oggetti che l’avversario ha indicato come utili, aumentando così le probabilità di un accordo reciproco.'}\\ \tt ARGOMENTAZIONE: {'Il set proposto fornisce il massimo valore di importanza per me (4184) rimanendo entro il limite di sforzo. Include tutti gli oggetti che avete mostrato interesse (C00, B35, B86, B32) e aggiunge gli oggetti ad alto valore per me (A77, B33, C54), garantendo un vantaggio significativo a entrambe le parti.'}\\ \tt PROPOSTA: {'C00', 'B35', 'A77', 'B86', 'B32', 'C54', 'B33'} 
	  } 
	   } 
	   } 
	  \\ 
 

    \theutterance \stepcounter{utterance}  

    & & \multicolumn{4}{p{0.6\linewidth}}{\cellcolor[rgb]{0.9,0.9,0.9}{%
	\makecell[{{p{\linewidth}}}]{% 
	  \tt {\tiny [A$\langle$GM]}  
	 ARGOMENTAZIONE: {'Il set proposto fornisce il massimo valore di importanza per me (4184) rimanendo entro il limite di sforzo. Include tutti gli oggetti che avete mostrato interesse (C00, B35, B86, B32) e aggiunge gli oggetti ad alto valore per me (A77, B33, C54), garantendo un vantaggio significativo a entrambe le parti.'}\\ \tt PROPOSTA: {'C00', 'B35', 'A77', 'B86', 'B32', 'C54', 'B33'} 
	  } 
	   } 
	   } 
	 & \\ 
 

    \theutterance \stepcounter{utterance}  

    & \multicolumn{4}{p{0.6\linewidth}}{\cellcolor[rgb]{0.8,1,0.9}{%
	\makecell[{{p{\linewidth}}}]{% 
	  \tt {\tiny [A$\rangle$GM]}  
	 RAGIONAMENTO STRATEGICO: {'Ho calcolato l\'importanza che otterrei accettando la vostra proposta (C00, B35, A77, B86, B32, C54, B33): solo 1 610 punti, ben al di sotto del valore che posso raggiungere mantenendo tutti gli oggetti che voi ritenete utili (C00, B35, B86, B32) e aggiungendo gli oggetti ad altissima importanza per me (A71, A04, A87, B41, C01, C54). Il mio set migliore è {'C00','A71','B35','A04','C01','C54','A87','B41','B86','B32'} con sforzo totale 3445 ≤ 3481 e importanza 4 690, il massimo possibile dato il vincolo. Perciò rifiuto la vostra offerta e propongo il mio set ottimale.'}\\ \tt RIFIUTO: {'C00', 'B35', 'A77', 'B86', 'B32', 'C54', 'B33'}\\ \tt ARGOMENTAZIONE: {'La vostra proposta fornisce un valore di importanza molto basso per me rispetto al set che propongo, che massimizza il mio punteggio mantenendo tutti gli oggetti chiave desiderati da entrambi e rimanendo entro il limite di sforzo.'}\\ \tt PROPOSTA: {'C00', 'A71', 'B35', 'A04', 'C01', 'C54', 'A87', 'B41', 'B86', 'B32'} 
	  } 
	   } 
	   } 
	 & & \\ 
 

    \theutterance \stepcounter{utterance}  

    & & \multicolumn{4}{p{0.6\linewidth}}{\cellcolor[rgb]{0.9,0.9,0.9}{%
	\makecell[{{p{\linewidth}}}]{% 
	  \tt {\tiny [GM$\rangle$B]}  
	 RIFIUTO: {'C00', 'B35', 'A77', 'B86', 'B32', 'C54', 'B33'}\\ \tt ARGOMENTAZIONE: {'La vostra proposta fornisce un valore di importanza molto basso per me rispetto al set che propongo, che massimizza il mio punteggio mantenendo tutti gli oggetti chiave desiderati da entrambi e rimanendo entro il limite di sforzo.'}\\ \tt PROPOSTA: {'C00', 'A71', 'B35', 'A04', 'C01', 'C54', 'A87', 'B41', 'B86', 'B32'} 
	  } 
	   } 
	   } 
	 & \\ 
 

    \theutterance \stepcounter{utterance}  

    & & & \multicolumn{4}{p{0.6\linewidth}}{\cellcolor[rgb]{1,0.85,0.72}{%
	\makecell[{{p{\linewidth}}}]{% 
	  \tt {\tiny [GM$\langle$B]}  
	 RAGIONAMENTO STRATEGICO: {'Ho calcolato il valore totale di importanza per me per varie combinazioni rispettando il limite di sforzo 3481. La combinazione che massimizza il mio punteggio include tutti gli oggetti ad alto valore per me (C00, B35, B86, B32, C01, C54, B33, A77) e aggiunge A04, che è anche fra gli oggetti che l\'avversario ha indicato come importanti. Questa proposta utilizza 3273 di sforzo, ben al di sotto del limite, e raggiunge un valore di importanza 4670, il più alto possibile mantenendo gli oggetti chiave per l\'altro partecipante.'}\\ \tt ARGOMENTAZIONE: {'Propongo questo insieme perché fornisce il massimo valore di importanza per me (4670) restando entro il limite di sforzo, includendo tutti gli oggetti che avete indicato come importanti (C00, B35, B86, B32) e aggiungendo A04, che appare anche nella vostra lista di oggetti chiave.'}\\ \tt PROPOSTA: {'C00', 'B35', 'B86', 'B32', 'C01', 'C54', 'B33', 'A77', 'A04'} 
	  } 
	   } 
	   } 
	  \\ 
 

    \theutterance \stepcounter{utterance}  

    & & \multicolumn{4}{p{0.6\linewidth}}{\cellcolor[rgb]{0.9,0.9,0.9}{%
	\makecell[{{p{\linewidth}}}]{% 
	  \tt {\tiny [A$\langle$GM]}  
	 ARGOMENTAZIONE: {'Propongo questo insieme perché fornisce il massimo valore di importanza per me (4670) restando entro il limite di sforzo, includendo tutti gli oggetti che avete indicato come importanti (C00, B35, B86, B32) e aggiungendo A04, che appare anche nella vostra lista di oggetti chiave.'}\\ \tt PROPOSTA: {'C00', 'B35', 'B86', 'B32', 'C01', 'C54', 'B33', 'A77', 'A04'} 
	  } 
	   } 
	   } 
	 & \\ 
 

    \theutterance \stepcounter{utterance}  

    & \multicolumn{4}{p{0.6\linewidth}}{\cellcolor[rgb]{0.8,1,0.9}{%
	\makecell[{{p{\linewidth}}}]{% 
	  \tt {\tiny [A$\rangle$GM]}  
	 RAGIONAMENTO STRATEGICO: {'Ho calcolato l\'importanza totale che otterrei accettando la vostra proposta (C00, B35, B86, B32, C01, C54, B33, A77, A04): 2 874 punti con sforzo 3 140. Ho poi determinato il set possibile con il massimo valore di importanza per me entro il limite 3 481, includendo tutti gli oggetti ad alta importanza per me (A71, A04, A87, B41, C01, C54, B35, B86, B32, C00). Questo set rispetta lo sforzo totale (3 312 ≤3 481) e mi fornisce 4 690 punti di importanza, che è il massimo raggiungibile. Pertanto rifiuto la vostra offerta e presento il mio set ottimale.'}\\ \tt RIFIUTO: {'C00', 'B35', 'B86', 'B32', 'C01', 'C54', 'B33', 'A77', 'A04'}\\ \tt ARGOMENTAZIONE: {'La vostra proposta fornisce un valore di importanza molto inferiore rispetto al set che posso ottenere mantenendo tutti gli oggetti chiave per voi e aggiungendo i miei oggetti di alto valore. Il mio set rispetta il vincolo di sforzo e massimizza il mio punteggio personale.'}\\ \tt PROPOSTA: {'C00', 'A71', 'A04', 'A87', 'B41', 'C01', 'C54', 'B35', 'B86', 'B32'} 
	  } 
	   } 
	   } 
	 & & \\ 
 

    \theutterance \stepcounter{utterance}  

    & & \multicolumn{4}{p{0.6\linewidth}}{\cellcolor[rgb]{0.9,0.9,0.9}{%
	\makecell[{{p{\linewidth}}}]{% 
	  \tt {\tiny [GM$\rangle$B]}  
	 RIFIUTO: {'C00', 'B35', 'B86', 'B32', 'C01', 'C54', 'B33', 'A77', 'A04'}\\ \tt ARGOMENTAZIONE: {'La vostra proposta fornisce un valore di importanza molto inferiore rispetto al set che posso ottenere mantenendo tutti gli oggetti chiave per voi e aggiungendo i miei oggetti di alto valore. Il mio set rispetta il vincolo di sforzo e massimizza il mio punteggio personale.'}\\ \tt PROPOSTA: {'C00', 'A71', 'A04', 'A87', 'B41', 'C01', 'C54', 'B35', 'B86', 'B32'} 
	  } 
	   } 
	   } 
	 & \\ 
 

    \theutterance \stepcounter{utterance}  

    & & & \multicolumn{4}{p{0.6\linewidth}}{\cellcolor[rgb]{1,0.85,0.72}{%
	\makecell[{{p{\linewidth}}}]{% 
	  \tt {\tiny [GM$\langle$B]}  
	 RAGIONAMENTO STRATEGICO: {'Ho calcolato l\'importanza totale per me delle varie combinazioni possibili rispettando il vincolo di sforzo ≤ 3481. Il set più vantaggioso include tutti gli oggetti ad alto valore per me (C00, B35, B86, B32, C01, C54) e aggiunge A77 (imp 668, effort 453) e C17 (imp 583, effort 777). Per completare il budget ho aggiunto B41 (imp 138, effort 542), che è accettabile per l’avversario e non penalizza significativamente il valore totale. Questo insieme utilizza 3329 di sforzo (← ≤ 3481) e fornisce un valore di importanza per me pari a 4514, sostanzialmente superiore alla proposta corrente dell’avversario (3546). Inoltre mantiene tutti gli oggetti chiave che l’avversario ha finora indicato come importanti (C00, B35, B86, B32, C01, C54, B41), aumentando la probabilità di un accordo reciprocamente vantaggioso.'} \\ \tt ARGOMENTAZIONE: {'Il set proposto massimizza il mio punteggio di importanza (4514) senza superare il limite di sforzo, mantenendo tutti gli oggetti che l’avversario ha già mostrato interesse (C00, B35, B86, B32, C01, C54, B41) e aggiungendo A77 e C17, i quali hanno elevata importanza per me.'} \\ \tt PROPOSTA: {'C00', 'C01', 'C54', 'B35', 'B86', 'B32', 'A77', 'C17', 'B41'} 
	  } 
	   } 
	   } 
	  \\ 
 

    \theutterance \stepcounter{utterance}  

    & & \multicolumn{4}{p{0.6\linewidth}}{\cellcolor[rgb]{0.9,0.9,0.9}{%
	\makecell[{{p{\linewidth}}}]{% 
	  \tt {\tiny [A$\langle$GM]}  
	 ARGOMENTAZIONE: {'Il set proposto massimizza il mio punteggio di importanza (4514) senza superare il limite di sforzo, mantenendo tutti gli oggetti che l’avversario ha già mostrato interesse (C00, B35, B86, B32, C01, C54, B41) e aggiungendo A77 e C17, i quali hanno elevata importanza per me.'} \\ \tt PROPOSTA: {'C00', 'C01', 'C54', 'B35', 'B86', 'B32', 'A77', 'C17', 'B41'} 
	  } 
	   } 
	   } 
	 & \\ 
 

    \theutterance \stepcounter{utterance}  

    & \multicolumn{4}{p{0.6\linewidth}}{\cellcolor[rgb]{0.8,1,0.9}{%
	\makecell[{{p{\linewidth}}}]{% 
	  \tt {\tiny [A$\rangle$GM]}  
	 RAGIONAMENTO STRATEGICO: {'Ho calcolato l\'importanza che otterrei accettando la vostra ultima proposta (C00, C01, C54, B35, B86, B32, A77, C17, B41) ed è solo 2718, molto inferiore al massimo che posso raggiungere mantenendo tutti gli oggetti che voi avete già indicato come importanti. Propongo quindi di rifiutare la vostra offerta e di proporre un set che conserva tutti i vostri oggetti chiave (C00, B35, B86, B32, C01, C54, B41) e aggiunge gli oggetti ad altissima importanza per me (A71, A04, A87), rimanendo entro il limite di sforzo (3445 ≤ 3481).'}\\ \tt RIFIUTO: {'C00', 'C01', 'C54', 'B35', 'B86', 'B32', 'A77', 'C17', 'B41'}\\ \tt ARGOMENTAZIONE: {'La vostra proposta non include tutti gli oggetti che per voi sono importanti (B41, B86, B32, ecc.) e, soprattutto, mi fornisce un valore di importanza molto basso. Il set che propongo mantiene tutti i vostri oggetti chiave e aggiunge A71, A04 e A87, che hanno un valore molto elevato per me, garantendo al contempo che lo sforzo totale sia entro il limite condiviso.'}\\ \tt PROPOSTA: {'C00', 'A71', 'B35', 'A04', 'C01', 'C54', 'A87', 'B41', 'B86', 'B32'} 
	  } 
	   } 
	   } 
	 & & \\ 
 

    \theutterance \stepcounter{utterance}  

    & & \multicolumn{4}{p{0.6\linewidth}}{\cellcolor[rgb]{0.9,0.9,0.9}{%
	\makecell[{{p{\linewidth}}}]{% 
	  \tt {\tiny [GM$\rangle$B]}  
	 RIFIUTO: {'C00', 'C01', 'C54', 'B35', 'B86', 'B32', 'A77', 'C17', 'B41'}\\ \tt ARGOMENTAZIONE: {'La vostra proposta non include tutti gli oggetti che per voi sono importanti (B41, B86, B32, ecc.) e, soprattutto, mi fornisce un valore di importanza molto basso. Il set che propongo mantiene tutti i vostri oggetti chiave e aggiunge A71, A04 e A87, che hanno un valore molto elevato per me, garantendo al contempo che lo sforzo totale sia entro il limite condiviso.'}\\ \tt PROPOSTA: {'C00', 'A71', 'B35', 'A04', 'C01', 'C54', 'A87', 'B41', 'B86', 'B32'} 
	  } 
	   } 
	   } 
	 & \\ 
 

    \theutterance \stepcounter{utterance}  

    & & & \multicolumn{4}{p{0.6\linewidth}}{\cellcolor[rgb]{1,0.85,0.72}{%
	\makecell[{{p{\linewidth}}}]{% 
	  \tt {\tiny [GM$\langle$B]}  
	 RAGIONAMENTO STRATEGICO: {'Ho verificato che la proposta dell’avversario supera il limite di sforzo (4221 > 3481). Per rispondere, ho calcolato il sotto‑insieme più vicino che rispetti il vincolo mantenendo tutti gli oggetti che sembrano importanti per l’avversario (C00, B35, B86, B32, C01, C54) e includendo anche A04 e B41, anch’essi richiesti dall’avversario. Questo sotto‑insieme ha uno sforzo totale di 3336, ben entro il limite, e conserva gran parte degli oggetti chiave della loro proposta, aumentando la probabilità di trovare un accordo reciproco.'} \\ \tt ARGOMENTAZIONE: {'Il set proposto rispetta il vincolo di sforzo (3336 ≤ 3481) e contiene tutti gli oggetti che sono stati indicati come importanti dall’avversario (C00, B35, B86, B32, C01, C54) oltre a includere A04 e B41, i quali erano esplicitamente richiesti nella loro ultima proposta. Questo compromesso massimizza il valore per entrambe le parti evitando di superare il limite condiviso.'} \\ \tt PROPOSTA: {'C00', 'B35', 'B86', 'B32', 'C01', 'C54', 'A04', 'B41'} 
	  } 
	   } 
	   } 
	  \\ 
 

    \theutterance \stepcounter{utterance}  

    & & \multicolumn{4}{p{0.6\linewidth}}{\cellcolor[rgb]{0.9,0.9,0.9}{%
	\makecell[{{p{\linewidth}}}]{% 
	  \tt {\tiny [A$\langle$GM]}  
	 ARGOMENTAZIONE: {'Il set proposto rispetta il vincolo di sforzo (3336 ≤ 3481) e contiene tutti gli oggetti che sono stati indicati come importanti dall’avversario (C00, B35, B86, B32, C01, C54) oltre a includere A04 e B41, i quali erano esplicitamente richiesti nella loro ultima proposta. Questo compromesso massimizza il valore per entrambe le parti evitando di superare il limite condiviso.'} \\ \tt PROPOSTA: {'C00', 'B35', 'B86', 'B32', 'C01', 'C54', 'A04', 'B41'} 
	  } 
	   } 
	   } 
	 & \\ 
 

    \theutterance \stepcounter{utterance}  

    & \multicolumn{4}{p{0.6\linewidth}}{\cellcolor[rgb]{0.8,1,0.9}{%
	\makecell[{{p{\linewidth}}}]{% 
	  \tt {\tiny [A$\rangle$GM]}  
	 RAGIONAMENTO STRATEGICO: {'Ho valutato la proposta corrente dell’avversario (C00, B35, B86, B32, C01, C54, A04, B41) che mi garantisce 3239 punti di importanza con sforzo 2560. Con il margine di sforzo disponibile (3481‑2560=922) posso includere gli oggetti ad altissima importanza per me, A71 (imp 783, s 332) e A87 (imp 668, s 553). L\'insieme risultante rispetta il limite di sforzo (3445 ≤ 3481) e porta il mio punteggio a 4690, il valore più alto possibile data la capacità. Inoltre questi due oggetti hanno anche un valore elevato per l’avversario, rendendo la proposta vantaggiosa per entrambe le parti.'}\\ \tt ARGOMENTAZIONE: {'Propongo di aggiungere A71 e A87 al vostro set corrente, così otteniamo un insieme che massimizza il mio valore di importanza (4690) rispettando il limite di sforzo, e allo stesso tempo fornisce un alto valore anche a voi (A71 783, A87 668).'}\\ \tt PROPOSTA: {'C00', 'B35', 'B86', 'B32', 'C01', 'C54', 'A04', 'B41', 'A71', 'A87'} 
	  } 
	   } 
	   } 
	 & & \\ 
 

    \theutterance \stepcounter{utterance}  

    & & \multicolumn{4}{p{0.6\linewidth}}{\cellcolor[rgb]{0.9,0.9,0.9}{%
	\makecell[{{p{\linewidth}}}]{% 
	  \tt {\tiny [GM$\rangle$B]}  
	 ARGOMENTAZIONE: {'Propongo di aggiungere A71 e A87 al vostro set corrente, così otteniamo un insieme che massimizza il mio valore di importanza (4690) rispettando il limite di sforzo, e allo stesso tempo fornisce un alto valore anche a voi (A71 783, A87 668).'}\\ \tt PROPOSTA: {'C00', 'B35', 'B86', 'B32', 'C01', 'C54', 'A04', 'B41', 'A71', 'A87'} 
	  } 
	   } 
	   } 
	 & \\ 
 

    \theutterance \stepcounter{utterance}  

    & & & \multicolumn{4}{p{0.6\linewidth}}{\cellcolor[rgb]{1,0.85,0.72}{%
	\makecell[{{p{\linewidth}}}]{% 
	  \tt {\tiny [GM$\langle$B]}  
	 RAGIONAMENTO STRATEGICO: {'Ho valutato la vostra proposta corrente, il suo sforzo totale è 3445 ≤ 3481, e il mio valore di importanza è 3546, un risultato accettabile rispetto a possibili alternative. Accettare garantisce un accordo rapido e soddisfa entrambi, includendo gli oggetti ad alto valore per voi (A71, A87) e mantenendo tutti gli oggetti chiave per me.'}\\ \tt ARGOMENTAZIONE: {'Accetto la vostra proposta perché rispetta il limite di sforzo, fornisce un valore di importanza significativo per me (3546) e include gli oggetti più importanti per voi, facilitando un accordo equo.'}\\ \tt ACCORDO: {'C00', 'B35', 'B86', 'B32', 'C01', 'C54', 'A04', 'B41', 'A71', 'A87'} 
	  } 
	   } 
	   } 
	  \\ 
 

    \theutterance \stepcounter{utterance}  

    & & & \multicolumn{2}{p{0.3\linewidth}}{\cellcolor[rgb]{0.95,0.95,0.95}{%
	\makecell[{{p{\linewidth}}}]{% 
	  \tt {\tiny [GM$|$GM]}  
	 game successful 
	  } 
	   } 
	   } 
	 & & \\ 
 

    \theutterance \stepcounter{utterance}  

    & & & \multicolumn{2}{p{0.3\linewidth}}{\cellcolor[rgb]{0.95,0.95,0.95}{%
	\makecell[{{p{\linewidth}}}]{% 
	  \tt {\tiny [GM$|$GM]}  
	 end game 
	  } 
	   } 
	   } 
	 & & \\ 
 

\end{supertabular}
}

\end{document}
